% 金星（综述）
% license CCBYSA3
% type Wiki

本文根据 CC-BY-SA 协议转载翻译自维基百科\href{https://en.wikipedia.org/wiki/Venus}{相关文章}。

金星是距离太阳第二近的行星。由于其轨道最接近地球，且两者同为类地岩石行星，并在大小、质量和表面引力上几乎完全一致，因此金星常被称为地球的“孪生”或“姐妹”行星。然而，金星在许多方面与地球显著不同，尤其是它没有液态水，其大气层远比太阳系中其他任何岩石天体都更加厚重和致密。金星大气主要由二氧化碳组成，并覆盖着一层浓密的硫酸云层，环绕整个行星。在平均地表高度处，大气温度达到 737\,K（464\,°C；867\,°F），压力为地球海平面的 92 倍，使其最低层大气呈现超临界流体状态。从地球观测时，金星呈现为类似恒星的亮点，比天空中任何其他天然光源都更明亮\(^\text{[22][23]}\)，并且作为一颗内行星，它总是出现在接近太阳的方向，要么作为最明亮的“启明星”，要么作为“长庚星”。

金星与地球的轨道使两者在约 1.6 年的会合周期中彼此接近。在此过程中，金星会比其他任何行星更接近地球；相较之下，由于水星的轨道更靠近太阳，因此它在自身轨道路径中整体上比其他行星更接近地球。在从地球出发的行星际空间飞行中，金星常被用作引力助推的中转点，从而提供更快速且更经济的航程。金星没有卫星，并以非常缓慢的逆行方式绕其自转轴旋转，这被认为是太阳潮汐锁定效应与金星巨大大气层的受热差异相互竞争的结果。因此，金星的一昼夜长达 116.75 个地球日，而其一个太阳年为 224.7 个地球日。

金星具有微弱的磁层；由于缺乏内部发电机，其磁场是由太阳风与大气相互作用所诱导产生的。在内部结构上，金星具有核心、地幔和地壳。内部热量通过活跃的火山活动逸出\(^\text{[24][25]}\)，从而导致地表更新，而不是板块构造。金星在其早期历史中可能曾拥有液态地表水及宜居环境\(^\text{[26][27]}\)，但由于失控的温室效应，所有水分最终蒸发，使金星演变为当前的状态\(^\text{[28][29][30]}\)。在云层高度处的大气条件是整个太阳系中与地球最为相似的，这些条件被认为可能有利于金星生命的存在，并在 2020 年发现了潜在的生物标志物，从而推动了新的研究与探测任务。

在人类历史上，全世界的人们都曾观测过金星，并使其在众多文化中具有特殊的重要性。随着望远镜的出现，金星的相位变得可辨识，并在 1613 年被呈现为推翻当时占主导地位的地心模型、支持日心模型的决定性证据。1961 年，“金星 1 号”首次造访金星，飞掠行星并实现了人类第一次行星际飞行。1962 年，第二次行星际任务“水手 2 号”带回了首批来自金星的数据。1967 年，首个行星际撞击器“金星 4 号”抵达金星，随后在 1970 年由着陆器“金星 7 号”成功登陆。截至 2025 年，“太阳轨道飞行器”正前往于 2026 年飞掠金星，而下一项计划发射前往金星的任务是“金星生命探测器”，预定同样于 2026 年发射。
\subsection{物理特征}
\begin{figure}[ht]
\centering
\includegraphics[width=8cm]{./figures/23fb842f59355847.png}
\caption{金星（从左数第二，伪彩色显示）按比例与太阳系内侧具有行星质量的天体一起展示，排列顺序依其轨道自太阳向外依次排列（从左到右：水星、金星、地球、月球、火星和谷神星）。} \label{fig_Venus_1}
\end{figure}
\begin{figure}[ht]
\centering
\includegraphics[width=8cm]{./figures/f163eb810133acdf.png}
\caption{金星在不同波长下的成像} \label{fig_Venus_2}
\end{figure}
金星是太阳系中四颗类地行星之一，这意味着它与地球一样属于岩石天体。它在大小和质量上与地球相似，因此常被描述为地球的“姐妹”或“孪生”行星\(^\text{[31]}\)。由于其自转缓慢，金星的形状非常接近球形\(^\text{[32]}\)。其直径为 12,103.6\,km（7,520.8\,mi），仅比地球小 638.4\,km（396.7\,mi）；其质量为地球的 81.5\%，使金星成为太阳系中第三小的行星。金星表面的环境与地球完全不同，因为其致密大气中 96.5\% 为二氧化碳，导致强烈的温室效应，其余约 3.5\% 为氮气\(^\text{[33]}\)。金星表面气压为 9.3\,MPa（93\,bar），平均表面温度为 737\,K（464\,°C；867\,°F），高于两种主要组分的临界点，使其地表大气成为一种超临界流体，主要由超临界二氧化碳及部分超临界氮组成。
\subsubsection{自然历史}
\textbf{形成}

包括金星在内的类地岩石行星被认为经历了五个阶段形成：尘埃沉降、微行星形成、行星胚胎阶段、巨撞阶段，以及最终的大气形成阶段。来自金星的有限测量数据使其形成时间线难以进行更详细的分析\(^\text{[34]}\)。

\textbf{未来}

金星预计将在七至八十亿年后当太阳演变为红巨星时，与水星一起被毁灭，也有可能包括地球和月球\(^\text{[35]}\)。
\subsubsection{地理}
\begin{figure}[ht]
\centering
\includegraphics[width=8cm]{./figures/884bdec4a9f38b8d.png}
\caption{彩色编码的高程地图，显示金星上呈黄色的高地“大陆”以及次要地形特征} \label{fig_Venus_3}
\end{figure}
\begin{figure}[ht]
\centering
\includegraphics[width=8cm]{./figures/4ef3c5c0b6c896cf.png}
\caption{金星表面雷达数据的球形视图，突出显示地表特征（1989 年，麦哲伦号）。这些颜色并不代表地表的真实外观[36]。} \label{fig_Venus_4}
\end{figure}
在 20 世纪探测器揭示部分秘密之前，金星表面一直是推测的对象。1975 年和 1982 年的“金星号”着陆器传回了由沉积物和相对棱角分明的岩石覆盖的地表图像\(^\text{[37]}\)。1990–91 年间，“麦哲伦号”对其表面进行了详细测绘。有大量火山活动的证据，并且大气中二氧化硫的变化可能表明存在活火山\(^\text{[38][39]}\)。

大约 80\% 的金星表面被光滑的火山平原覆盖，其中 70\% 为带有皱褶脊的平原，10\% 为光滑或叶状平原\(^\text{[40]}\)。其余的地表由两个高地“大陆”构成，一个位于行星的北半球，另一个位于赤道以南。北部高地称为伊什塔尔高原（Ishtar Terra），得名于巴比伦的爱神伊什塔尔，其面积约与澳大利亚相当。麦克斯韦山脉位于伊什塔尔高原，其最高峰斯卡迪山（Skadi Mons）是金星上的最高点，海拔比金星平均表面高度高 11\,km（7\,mi）\(^\text{[41]}\)。南部高地称为阿芙洛狄忒高原（Aphrodite Terra），得名于希腊神话中的爱神，是两大高地中较大的一个，面积约与南美洲相当。该区域的大部分被裂缝和断层网络覆盖\(^\text{[42]}\)。

近期有关于金星上岩浆流动的证据（2024 年）\(^\text{[43]}\)，例如在盾状火山 Sif Mons 的熔岩流以及在平坦平原 Niobe Planitia 上的熔岩流\(^\text{[24]}\)。地表可见破火山口。金星的撞击坑数量很少，表明其表面相对年轻，约为 3 亿至 6 亿年\(^\text{[44][46]}\)。金星除了拥有岩石行星上常见的撞击坑、山脉和山谷外，还具有一些独特的地表特征。其中包括称为“farra”的平顶火山构造，其形状类似煎饼，直径从 20 至 50\,km（12 至 31\,mi）不等，高度在 100 至 1,000\,m（330 至 3,280\,ft）之间；称为“novae”的放射状、星形裂缝系统；具有放射状和同心裂隙、类似蜘蛛网的构造，被称为“arachnoids”；以及“coronae”，即被洼地包围的环状裂缝结构。这些特征均具有火山成因\(^\text{[46]}\)。
\begin{figure}[ht]
\centering
\includegraphics[width=8cm]{./figures/5421d90f38272eab.png}
\caption{彩色化影像（“金星 9 号”，1975 年），金星天空在地表呈橙黄色，这是由于瑞利散射或低层大气中的蓝色吸收体所致，而在更高的高度则呈白色[47][48]；同时其地表为类似玄武岩的深灰色，可能因氧化而呈红色[36]。} \label{fig_Venus_5}
\end{figure}
大多数金星表面特征以历史和神话中的女性命名\(^\text{[49]}\)。例外包括麦克斯韦山脉，以詹姆斯·克拉克·麦克斯韦命名，以及阿尔法高地区（Alpha Regio）、贝塔高地区（Beta Regio）和奥夫达高地区（Ovda Regio）。这后三个特征是在国际天文学联合会（负责行星命名的机构）采用现行系统之前命名的\(^\text{[50]}\)。

金星地形特征的经度以其本初子午线为基准表达。最初的本初子午线通过位于阿尔法高地区南部、椭圆形地形“夏娃”（Eve）中央的雷达亮点\(^\text{[51]}\)。在“金星号”任务完成后，本初子午线被重新定义为通过塞德娜平原（Sedna Planitia）上阿里阿德涅（Ariadne）陨石坑中央峰的位置\(^\text{[9][52]}\)。

地层学上最古老的镶嵌地形（tessera terrains）在“金星快车号”和“麦哲伦号”测量中显示出相对于周围玄武质平原更低的热辐射率，这表明其具有不同的、可能更富长英质（felsic）的矿物组合\(^\text{[29][53]}\)。生成大量长英质地壳的机制通常需要存在海洋和板块构造，这意味着早期金星上可能存在宜居条件，并在某一时期拥有大量水体\(^\text{[54]}\)。然而，镶嵌地形的具体性质仍不确定\(^\text{[55]}\)。

2023 年的研究首次提出金星在古代可能存在板块构造，因此可能拥有更宜居的环境，甚至可能具备维持生命的能力\(^\text{[26][27]}\)。金星因此成为研究类地行星形成及其宜居性的重要案例。

\textbf{火山活动}
\begin{figure}[ht]
\centering
\includegraphics[width=8cm]{./figures/96e9067c358877bf.png}
\caption{金星埃斯特拉区域内两个煎饼穹丘的雷达拼接图——二者均宽 65\,km（40\,mi），高度不足 1\,km（0.62\,mi）} \label{fig_Venus_6}
\end{figure}
金星表面的很大一部分似乎是由火山活动塑造的。金星上的火山数量是地球的数倍，并拥有 167 座直径超过 100\,km（60\,mi）的大型火山。地球上唯一与其规模相当的火山构造是夏威夷大岛\(^\text{[46]:154}\)。金星上已识别并绘制地图的火山数量超过 85,000 座\(^\text{[56][57]}\)。这并不是因为金星的火山活动比地球更活跃，而是因为其地壳更古老，且没有经历地球上活跃的侵蚀过程。地球的海洋地壳在构造板块边界通过俯冲不断被循环，其平均年龄约为 1 亿年\(^\text{[58]}\)，而金星表面估计为 3 亿至 6 亿年\(^\text{[44][46]}\)。

多项证据指向金星上持续存在火山活动。上层大气中的二氧化硫浓度在 1978 年至 1986 年间下降了 10 倍，在 2006 年急剧上升，然后再次下降 10 倍\(^\text{[59]}\)。这可能意味着多次大型火山喷发提高了其含量\(^\text{[60][61]}\)。有人提出金星的闪电（见下文）可能源自火山活动（即火山闪电）。2020 年 1 月，天文学家报告发现证据表明金星目前存在火山活动，具体而言，他们检测到橄榄石这种火山产物，而该物质在金星表面会迅速风化\(^\text{[62][63]}\)。

这一巨大的火山活动由炽热的内部提供能量，根据模型，它可以由行星早期阶段的高能碰撞以及类似地球的放射性衰变来解释。撞击的速度会明显高于地球，一方面因为金星更靠近太阳而运动更快，另一方面因为高偏心率天体与金星相撞时具有更高速度\(^\text{[64]}\)。

在 2008 年和 2009 年，“金星快车号”首次观测到正在进行火山活动的直接证据，其形式是在裂谷区加尼斯峡谷（Ganis Chasma）内的四个短暂出现的局部红外热点\(^\text{[note 1]}\)，该区域靠近盾状火山玛特山（Maat Mons）\(^\text{[65]}\)。其中三个热点在连续轨道中被重复观测到。这些热点被认为代表新近喷发的熔岩\(^\text{[66][67]}\)。实际温度未知，因为热点的面积无法测量，但推测其温度可能在 800–1,100\,K（527–827\,°C；980–1,520\,°F）范围内，而正常温度约为 740\,K（467\,°C；872\,°F）\(^\text{[68]}\)。2023 年，科学家重新分析了“麦哲伦号”拍摄的玛特山区域地形图像。通过计算机模拟，他们确认该区域在 8 个月间发生了地形变化，并得出其原因是活跃火山活动\(^\text{[69]}\)。

\textbf{陨石坑}
\begin{figure}[ht]
\centering
\includegraphics[width=8cm]{./figures/d8116d164677ab95.png}
\caption{金星表面的撞击坑（伪彩色，由雷达数据重建的三维投影影像）} \label{fig_Venus_7}
\end{figure}
金星上有将近一千个撞击坑，均匀分布在其表面。在其他具有撞击坑的天体上，例如地球和月球，撞击坑呈现出不同程度的退化状态。月球上的退化由后续撞击造成，而地球上的退化则由风化和降雨侵蚀导致。在金星上，约 85\% 的撞击坑保持完好状态。撞击坑的数量及其良好的保存状况表明，该行星曾在 3 亿至 6 亿年前经历过一次全球性地表重塑事件\(^\text{[44][45]}\)，随后火山活动逐渐衰减\(^\text{[70]}\)。地球的地壳处于持续运动中，而金星被认为无法维持此类过程。在没有板块构造来散发地幔热量的情况下，金星经历一种循环过程，即地幔温度不断升高，直到达到削弱地壳的临界水平。随后，在大约 1 亿年的时间内，发生大规模的俯冲，将地壳完全循环\(^\text{[46]}\)。

金星的撞击坑直径范围从 3 到 280\,km（2 到 174\,mi）。没有直径小于 3\,km 的撞击坑，这是由于致密大气对进入物体的影响所致。动能低于某一阈值的物体会被大气大幅减速，以至于无法形成撞击坑\(^\text{[71]}\)。直径小于 50\,m（160\,ft）的入射物体将在到达地面之前于大气层中解体并焚毁\(^\text{[72]}\)。
\subsubsection{内部结构}
\begin{figure}[ht]
\centering
\includegraphics[width=8cm]{./figures/dbfa21d636bc841c.png}
\caption{金星的分化结构} \label{fig_Venus_8}
\end{figure}
由于缺乏反射地震学的数据或其转动惯量的测量结果，人们对金星内部结构和地球化学性质的直接信息非常有限\(^\text{[73]}\)。金星与地球在大小和密度上的相似性表明，它们可能拥有相似的内部结构：核心、地幔和地壳。与地球相似，金星的核心很可能至少部分为液态，因为这两颗行星的冷却速率大致相同\(^\text{[74]}\)，尽管完全固态核心也不能被排除\(^\text{[75]}\)。金星体积略小意味着其深部内部的压强比地球低 24\%\(^\text{[76]}\)。根据行星模型预测的转动惯量数值，金星核心半径估计为 2,900–3,450\,km\(^\text{[75]}\)。基于 2006 至 2020 年之间测量的自转轴进动速率所推算的转动惯量，目前的估计为 3,500\,km\(^\text{[13][77]}\)。

金星地壳平均估计厚度为 40\,km，最大可达 65\,km\(^\text{[78]}\)。

两颗行星的主要差异在于金星缺乏板块构造的证据，可能原因是其地壳太坚固，在没有水减少黏滞性的情况下无法发生俯冲。这导致行星热量损失较少，阻碍其冷却，并为其缺乏内部磁场提供了可能的解释\(^\text{[79]}\)。金星可能改以周期性的全球地表重塑事件释放内部热量\(^\text{[44]}\)。
\subsubsection{磁场与核心}
\begin{figure}[ht]
\centering
\includegraphics[width=8cm]{./figures/81a077950d438b05.png}
\caption{} \label{fig_Venus_9}
\end{figure}
1967 年，“金星 4 号”发现金星的磁场比地球弱得多。该磁场是由电离层与太阳风之间的相互作用所诱导产生的\(^\text{[80][81][page needed]}\)，而不是像地球那样由核心内部的发电机产生。金星这种微弱的诱导磁层对大气抵御太阳和宇宙辐射的保护几乎可以忽略不计。

金星缺乏本征磁场令科学家感到意外，因为金星在大小上与地球相似，原本应在其核心中存在发电机。发电机需要三个条件：导电液体、自转和对流。金星的核心被认为具有电导性，而尽管其自转通常被认为过于缓慢，模拟表明其仍足以产生发电机\(^\text{[82][83]}\)。这意味着缺乏发电机的原因在于金星核心内部缺乏对流。在地球上，对流发生在核心的液态外层中，因为液态层底部的温度远高于顶部。在金星上，一次全球性的地表重塑事件可能终止了板块构造，导致通过地壳的热通量减小。这种隔热效应会使地幔温度升高，从而减少核心的热通量。因此，无法提供用于驱动磁场的内部地球发电机，核心释放的热量则重新加热地壳\(^\text{[84]}\)。

一种可能性是金星没有固态内核\(^\text{[85]}\)，或者其核心没有在冷却，使得整个液态核心的温度大致相同。另一种可能性是其核心已经完全固化。核心的状态高度依赖于硫的浓度，而当前尚不清楚\(^\text{[84]}\)。

另一种可能性是金星缺乏像地球“形成月球”事件那样的大型撞击，使金星核心在增量形成过程中产生分层结构，而没有足够的力量触发或维持对流，从而无法形成“地质发电机”\(^\text{[86]}\)。

金星周围微弱的磁层意味着太阳风会直接与其外层大气相互作用。在这里，水分子因紫外线辐解而生成氢和氧离子。随后，太阳风提供能量，使其中一些离子获得足够速度逃离金星的引力场。该侵蚀过程导致轻质的氢、氦和氧离子持续流失，而高质量分子如二氧化碳更可能被保留。太阳风导致的大气侵蚀可能造成金星在形成后的前十亿年内失去了大部分水\(^\text{[87]}\)。然而，该行星在最初的 20–30 亿年内可能保有发电机，因此水的流失也可能发生得更晚\(^\text{[88]}\)。这种侵蚀使得大气中高质量氘与低质量氢的比值相比太阳系其他区域增加了 100 倍\(^\text{[89]}\)。
\subsection{大气与气候}
\begin{figure}[ht]
\centering
\includegraphics[width=8cm]{./figures/8e5406aad8f1ba0c.png}
\caption{通过紫外成像显现的金星大气云层结构} \label{fig_Venus_10}
\end{figure}
金星拥有致密的大气层，由 96.5\% 的二氧化碳、3.5\% 的氮气组成——两者在金星地表均以超临界流体形式存在，其密度为水的 6.5\%\(^\text{[90]}\)——并含有微量其他气体，包括二氧化硫\(^\text{[91]}\)。其大气质量为地球的 92 倍，而地表压力约为地球的 93 倍，相当于地球海洋表面下近 1\,km（5⁄8\,mi）处的压力。地表大气密度为 65\,kg/m³（4.1\,lb/cu\,ft），为水的 6.5\%\(^\text{[90]}\)，或为地球海平面 293\,K（20\,°C；68\,°F）空气密度的 50 倍。富含 CO₂ 的大气产生了太阳系中最强的温室效应，使其表面温度至少达到 735\,K（462\,°C；864\,°F）\(^\text{[92][93]}\)。这使得金星表面比水星更热\(^\text{[94]}\)，尽管金星与太阳的距离几乎是水星的两倍，因此其接收的太阳辐照度仅为水星的四分之一，即 2,600\,W/m²（约为地球的两倍）\(^\text{[5]}\)。由于其失控温室效应，金星被卡尔·萨根等科学家视为与地球气候变化相关的警示和研究对象\(^\text{[95]}\)。因此金星被称为温室行星\(^\text{[96]}\)，或处于温室炼狱中的行星\(^\text{[97]}\)。

金星大气中富含原初稀有气体，与地球大气相比更为丰富\(^\text{[98]}\)。这种富集表明金星在演化早期便与地球发生分歧。一个异常大的彗星撞击\(^\text{[99]}\)或从太阳星云吸积更大质量的原始大气\(^\text{[100]}\)被提出用于解释这种富集。然而，大气中放射性氩–40 含量较低，这种元素可视为地幔脱气的替代指标，这表明主要岩浆作用在早期便已终止\(^\text{[101][102]}\)。

研究表明，在数十亿年前，金星的大气可能与早期地球的大气非常相似，并且其地表可能存在大量液态水\(^\text{[103][104][105]}\)。经过约 6 亿年至数十亿年的时间\(^\text{[106]}\)，太阳光度上升及可能发生的大规模火山重塑导致原始水体蒸发\(^\text{[107]}\)。当大气中的温室气体（包括水）达到临界值后，便触发了失控温室效应\(^\text{[108]}\)。尽管金星地表环境已不再适合任何可能在此事件前形成的类地生命，但人们推测金星表面上方 50\,km（30\,mi）的高层云层中可能存在生命，因为该高度的环境是整个太阳系中最类似地球的\(^\text{[109]}\)，温度介于 303 至 353\,K（30 至 80\,°C；86 至 176\,°F）之间，气压和辐射水平与地球表面相当，但大气成分为酸性云和二氧化碳\(^\text{[110][111][112]}\)。更具体而言，在 48 至 59\,km 的高度范围内，温度和辐射条件适宜生命生存；在更低的高度水会蒸发，而在更高的高度紫外辐射将过强\(^\text{[113][114]}\)。2020 年 9 月，有研究报告在金星大气中探测到磷化氢的吸收线，而该物质在已知条件下无法通过非生物途径生成，因此引发了关于金星大气中可能存在现存生命的猜测\(^\text{[115][116]}\)。后续研究将这一光谱信号解释为二氧化硫\(^\text{[117]}\)，或发现实际上不存在该吸收线\(^\text{[118][119]}\)。
\begin{figure}[ht]
\centering
\includegraphics[width=8cm]{./figures/b77f963239de7576.png}
\caption{按高度绘制的大气剖面图（左侧刻度）：云层、温度变化（VIRA 粗线及底部刻度）、压力变化（右侧刻度）以及风速（PV 虚线及顶部刻度）} \label{fig_Venus_11}
\end{figure}
热惯性以及下层大气中由风输送热量的作用意味着，尽管金星自转缓慢，其面对太阳与背向太阳的半球之间的地表温度并没有显著差异。地表的风速很慢，仅为每小时数公里，但由于地表大气密度极高，它们对障碍物施加的力非常大，并能够在地表搬运尘埃和小石块。即便不考虑其高温、高压和缺氧环境，这种阻力本身就足以使人在地表难以行走\(^\text{[120]}\)。

在致密的 CO₂ 层之上，距地表 45 至 70\,km 的高度存在厚厚的云层\(^\text{[121]}\)，主要由硫酸组成，其形成过程是二氧化硫分子在紫外辐射催化下与水发生反应\(^\text{[122]}\)，从而生成硫酸水合物\(^\text{[123]}\)。此外，云层中含有约 1\% 的三氯化铁\(^\text{[124]}\)。云粒子的其他可能成分包括硫酸铁、氯化铝和五氧化二磷。不同高度的云层具有不同的成分和粒径分布[124]。这些云层像地球的厚云层一样，将落在其上的阳光约 70\% 反射回太空[126]，并由于覆盖整个行星，使得金星表面无法被可见光直接观测。永久性的云层意味着尽管金星比地球更接近太阳，其地表接收到的阳光却更少，只有 10\% 的入射阳光能到达地表[127]，使得其地表白天平均照度约为 14,000 勒克斯，与地球“阴天白昼”的照度相当[128]。

金星大气的自转速度远快于其固体本体的自转，这一现象被称为大气超自转[129]。这导致云顶处出现强劲的 300\,km/h（185\,mph）风，它们绕行星一周仅需约 4 天，相当于行星自转速度的 60 倍[130]，而地球上最强的风速仅为其自转速度的 10–20\%。

尽管金星在可见光下看起来缺乏表面特征，但在紫外波段却存在条带或条纹，其成因尚未确定。紫外吸收可能来源于一种由氧与硫组成的化合物 OSSO，该分子中两个硫原子之间具有双键，并存在“顺式”和“反式”两种构型，或可能来源于 S₂ 至 S₈ 的多硫化合物[131]。
\begin{figure}[ht]
\centering
\includegraphics[width=8cm]{./figures/d0e3e9800a1218c4.png}
\caption{} \label{fig_Venus_12}
\end{figure}
金星的地表实际上是等温的；它不仅在两个半球之间保持恒定温度，而且在赤道与两极之间也保持恒定温度[5][133]。金星的自转轴倾角极小——不到 3°，相比之下地球为 23°——这也使季节性温度变化极小[134]。海拔是影响金星温度为数不多的因素之一。金星最高点——麦克斯韦山脉中的斯卡迪山——因此是金星上最冷的地点，温度约为 655\,K（380\,°C；715\,°F），气压约为 4.5\,MPa（45\,bar）[132][135]。1995 年，“麦哲伦号”探测器在金星最高山峰顶部成像到一种高反射物质，被称为“金星雪”，与地球上的雪极为相似。这种物质可能通过类似于雪的过程形成，但温度高得多。由于其挥发性过高而无法在地表凝结，它以气态上升至更高、更冷的高度，在那里可能发生沉降。该物质的成分尚未确定，但推测范围从碲单质到硫化铅（方铅矿）不等[136]。

尽管金星没有季节变化，天文学家在 2019 年发现其大气对阳光吸收存在周期性变化，可能由悬浮在高层云中的不透明吸收粒子引起。这种变化导致金星纬向风速出现观测到的变化，并似乎与太阳 11 年黑子周期同步[137]。

金星大气中是否存在闪电长期以来存在争议[138]，自从苏联“金星号”探测器首次探测到疑似闪电脉冲后[139][140][141]。在 2006–07 年，“金星快车号”清晰探测到了哨声模式波，这是闪电的特征信号。其间歇性出现表明其与天气活动相关。根据这些测量，其闪电发生率至少为地球的一半[142]，然而其他仪器完全未能探测到闪电[138]。闪电的起源仍不明确，但可能来自云层或金星火山。

2007 年，“金星快车号”发现金星南极存在一个巨大的双极大气涡旋[143][144]。“金星快车号”在 2011 年发现金星高层大气中存在臭氧层[145]。2013 年，欧洲航天局科学家报告称金星的电离层向外流动，其方式类似于“在类似条件下彗星电离尾的外泄”[146][147]。

在 2015 年 12 月，以及较弱的 2016 年 4 月和 5 月，参与日本“晓号”任务的研究人员在金星大气中观测到弧形结构。这被认为是太阳系中最大、可能也是首个直接证据确认的驻波重力波[148][149][150]。

金星地表大气的颜色与声音[151]已被记录：天空在地表呈橙黄色，而在更高的高度则呈白色[48]。
\subsection{轨道与自转}









